\section*{Chapter 2: Background}
\setcounter{section}{2} \setcounter{subsection}{0}
\phantomsection \addcontentsline{toc}{section}{\textit{Chapter 2: Background}}



\subsection{From Process-centric Systems to Agent-mediated Computing}

Modern operating systems are traditionally organised around the process as the primary abstraction for scheduling and protection.
Classical designs emphasise address-space isolation, controlled access to files and devices, and well-defined inter-process communication \cite{tanenbaum2022}. 
These abstractions remain fundamental, but they are increasingly strained by agent-based workloads.

In emerging agent systems, execution is no longer driven solely by programs issuing system calls. 
Instead, an agent plans over semantically described tools, constructs structured requests, and performs actions on behalf of a user. 
This introduces a layer of intent between the user and execution. 
As a result, correctness depends not only on process isolation, but also on whether the semantics of a requested operation match what is actually executed.

Recent work reflects this shift. 
Systems such as generative agents demonstrate long-lived entities capable of maintaining memory and coordinating workflows \cite{park2023generative}, while tool-oriented LLM systems show that performance depends critically on structured access to external capabilities \cite{qin2024toolllm}.
More recent autonomous agents, such as Manus\cite{mei2024aios,liu2026agentos}, further illustrate that agents can plan, execute, and verify complex tasks through multi-agent decomposition. 
In these settings, the system must preserve consistency between tool descriptions, permissions, and runtime behaviour not directly captured by traditional OS abstractions.
That mismatch is the starting point for the present work, motivating the question of which consistency guarantees are better placed in the kernel rather than delegated to a user-space layer.



\subsection{AI Agent Operating Systems}

Motivated by these challenges, several systems have begun to position themselves explicitly as AI agent operating systems. 
For example, recent proposals advocate embedding LLM-centric services into a system-level kernel that manages scheduling, memory, and tool access for agents \cite{packer2024memgpt, zhang2024llmos, wu2024oscopilot}. 
More broadly, the LLM-as-OS paradigm conceptualises the model as a system kernel, with agents acting as applications and tools as system resources \cite{packer2024memgpt}. 
These designs aim to provide unified abstractions for agent execution, including context management, multi-agent coordination, and resource governance.

In parallel, practical systems have emerged that approximate this vision at the application layer. 
OpenClaw\cite{openclaw2026, openclaw_security1, openclaw_security2} represents a widely deployed open-source agent framework that operates as a persistent, self-hosted assistant capable of executing real-world tasks through LLM-driven planning and tool invocation.
It exposes capabilities through a gateway and modular “skills,” effectively acting as a user-space control layer for coordinating tool execution.
Similarly, commercial platforms such as Apple Intelligence\cite{apple2024pcc} integrate AI agents deeply into the operating environment, combining on-device models with system services to provide personalised, context-aware assistance across applications.

At the same time, industrial systems are beginning to integrate AI agents more tightly into operating environments. 
These developments suggest that agent execution is increasingly becoming an operating-system concern rather than a purely application-level feature.

Despite differences in implementation, these systems share a common trend: they introduce an additional control layer that mediates between user intent, model reasoning, and tool execution.
This layer increasingly resembles operating-system functionality, even when implemented entirely in user space.
\texttt{linux-mcp} takes a more conservative premise: rather than rethinking the operating system for agents, it asks which specific invariants from that emerging control layer benefit from kernel placement, and whether placing only those invariants in the kernel is sufficient to obtain meaningful guarantees.



\subsection{What MCP Provides}

The MCP provides a useful abstraction for structuring this interaction. 
It defines a client-host-server architecture in which tools, resources and prompts are exposed through a standardised interface \cite{mcp2025}.
The host coordinates sessions and mediates access, while servers implement capabilities behind a JSON-RPC protocol.

This separation simplifies interoperability and modularity. 
However, it also concentrates responsibility in the host: identity management, approval decisions, and consistency between tool descriptions and execution all depend on correct host behaviour.
MCP defines how components communicate, but does not specify how a local system should enforce these properties.
\texttt{linux-mcp} implements the enforcement layer that MCP leaves unspecified: it provides the local mechanism that turns the protocol's abstract trust model into a set of kernel-verified invariants.

\subsection{Limits of User-space Mediation}

In most existing systems, including OpenClaw architectures, this control logic resides entirely in a user-space gateway \cite{langchain2022, autogpt2023}.
This design enables rapid development and flexible integration, but introduces several structural weaknesses.
Because tool invocation depends on semantically rich manifests---combining natural-language descriptions, schemas and examples \cite{schick2023toolformer, patil2023gorilla, mcp2025}---there is an implicit assumption that these descriptions remain aligned with the underlying implementation.
When both interpretation and enforcement reside in a mutable user-space process, inconsistencies can arise \cite{liu2023prompt}.
A change in implementation or a weakened control path may result in executing a capability that no longer matches the approved semantics.
A similar issue arises with identity.
Agent systems operate over richer notions of identity than traditional processes, often binding requests to sessions or agents rather than PIDs \cite{yao2023react, autogpt2023, park2023generative}.
Without anchoring these identities in OS-visible credentials, approval decisions can be reused or replayed across contexts, weakening isolation
guarantees \cite{greshake2023prompt}.
Finally, durability becomes problematic when control-plane state is held only in memory \cite{packer2024memgpt}.
Approval decisions and pending requests may be lost on crash or restart, leaving no authoritative state that other components can consult.
This fragility is particularly relevant in long-running agent systems where execution spans multiple steps and decisions \cite{wang2023voyager, yang2024swe}.
These limitations are not unique to any single framework.
Rather, they reflect a broader pattern: current agent systems increasingly rely on control plane invariants \cite{anderson1972computer} that are difficult to enforce when confined to user space \cite{garfinkel2003terra, watson2010capsicum, ji2026seagent, shi2025progent}.

\subsection{Attack Classes Outside the Scope of Kernel Admission Control}
\label{sec:attack-scope}

Kernel-level admission control addresses a specific subset of attack classes, and it is worth being explicit about its boundaries before positioning the design.
Two dominant attack categories in AI agent systems fall outside the scope of this approach.

\textbf{Prompt injection.}
In a prompt injection attack, adversarial content embedded in a tool's output or in the environment manipulates the LLM planner into choosing an unintended tool or constructing a malicious payload \cite{greshake2023prompt, liu2023prompt}.
The attack operates entirely at the semantic layer, before or independently of any \texttt{tool:exec} call reaching the control plane.
Defending against prompt injection requires defenses at the model or application layer---input sanitisation, output monitoring, sandboxed execution contexts---none of which are provided by admission control at the kernel level.
\texttt{linux-mcp} does not address prompt injection; it assumes that the payload submitted to \texttt{tool:exec} reflects genuine user or planner intent.

\textbf{Supply-chain compromise of manifests.}
If an attacker modifies \texttt{tool-app/manifests/*.json} before \texttt{mcpd} starts, the kernel will faithfully enforce the attacker's tool definitions.
Protecting manifest integrity requires file-system-level mechanisms such as IMA (Integrity Measurement Architecture) or dm-verity, which operate at a different layer and are orthogonal to this work.
\texttt{linux-mcp} treats on-disk manifests as a trusted input, consistent with the assumption that the host file system is not compromised.

\subsection{Design Positioning}

The present work sits between two extremes.
On one end, AIOS-style proposals advocate for a full rethinking of the operating system around LLM-centric abstractions.
On the other, existing agent frameworks operate entirely in user space, relying on application-level logic to enforce correctness and security.

The approach taken here is narrower.
Instead of redefining the entire operating system, it focuses on a specific question: which control-plane invariants in MCP-style tool invocation benefit from kernel-level enforcement, given that the gateway and manifests are trusted inputs?


\begin{table}[t]
\centering
\small
\begin{tabular}{l p{4.8cm} p{5.2cm}}
\toprule
\textbf{Attack class} & \textbf{Representative attack} & \textbf{Defense in \texttt{linux-mcp}} \\
\midrule

Spoofing 
& Forged session identity or mismatched tool hash 
& Kernel-verified identity binding and semantic-hash consistency \\

Replay 
& Reuse of stale approval tickets or prior requests 
& Kernel-enforced approval lifecycle with binding and lifetime validation \\

Substitution 
& Executing a tool diverging from declared semantics 
& Verification of manifest-derived semantic identity \\

Escalation 
& Bypassing gateway checks via logic flaws 
& Mandatory kernel arbitration for all invocations \\

\bottomrule
\end{tabular}
\caption{Control-plane threats and defenses (in scope).}
\label{tab:threat_in_scope}
\end{table}



\subsection{Relation to Native Linux Security Mechanisms}

Any such design must coexist with existing Linux security mechanisms. Facilities such as seccomp, AppArmor, SELinux and Landlock provide strong guarantees over system calls, resources and process privileges \cite{seccomp2026,lsm2026,apparmor2026,selinux2026,landlock2026}. 
These mechanisms operate at the level of kernel objects and access rights.

However, they do not capture higher-level concepts such as tool identity, semantic equivalence between a manifest and an execution endpoint, or the scope of an approval decision. 
For example, they cannot directly express whether a particular tool invocation matches an approved semantic description, or whether a request belongs to a specific agent session.

This gap becomes more pronounced as agent systems evolve toward persistent, autonomous operation with rich tool ecosystems.
Addressing it requires introducing a control layer that reasons about semantics and identity before execution reaches user-space logic.



\subsection{Related Work}
\label{sec:related-work}

Several existing mechanisms address related but distinct problems; understanding why they fall short for MCP-style enforcement motivates the present design.

\textbf{Seccomp-BPF.}
Seccomp-BPF filters system-call arguments at the process boundary \cite{seccomp2026}.
It can restrict which syscalls a gateway process may issue, but it operates at the level of kernel objects rather than tool identity: it cannot distinguish a legitimate tool invocation from a forged one carrying the same syscall pattern.
Applied to \texttt{mcpd}, seccomp can reduce the syscall surface but cannot enforce that the semantic hash in a request matches what was registered, nor that an approval ticket is unconsumed.
That gap is what the seccomp baseline in Chapter~\ref{} exposes: adding stricter syscall filtering reduces the bypass rate but does not eliminate it for attack classes that do not require a privileged syscall.

\textbf{eBPF and LSM hooks.}
eBPF programs can enforce byte-level socket policies with very low overhead \cite{lsm2026}, and LSM hooks attach to VFS and socket operations at well-defined kernel hook points \cite{krohn2007flume}.
Neither, however, can natively maintain a dynamic, per-tool hash registry and perform variable-length hash comparisons without significant complexity: the BPF verifier's loop and memory restrictions make runtime registry lookups impractical, and LSM hooks fire at the system-call boundary rather than at the application-protocol boundary where tool identity is meaningful (see Section~\ref{sec:design-alternatives} for a more detailed analysis).

\textbf{Landlock.}
Landlock provides process-scoped filesystem sandboxing through unprivileged LSM restrictions \cite{landlock2026}.
It is complementary to this work---a tool service could be Landlock-sandboxed to limit its filesystem reach---but it does not address tool identity, approval binding, or approval-state durability across daemon failures.

\textbf{Capsicum.}
Capsicum introduces a capability model in which each open file descriptor represents a scoped authority \cite{watson2010capsicum}.
Adapting this to MCP-style tool invocation would require tool services to be written against a capability API, violating the design goal that tool authors need not be aware of any kernel interface.
More fundamentally, Capsicum capabilities are per-file-descriptor and cannot natively encode whether a specific tool invocation matches an approved semantic description.

\textbf{AIOS and LLM-as-OS proposals.}
Several recent systems propose embedding LLM-specific abstractions into a new OS substrate \cite{zhang2024llmos, wu2024oscopilot, packer2024memgpt}.
These differ from this work in scope: they aim to replace or fundamentally extend the operating system for agent workloads.
\texttt{linux-mcp} takes the opposite approach---leaving the existing OS unchanged and adding only a narrow kernel module for the invariants that user space cannot reliably hold.

\textbf{MAC frameworks for LLM agents.}
Recent work has begun applying mandatory access control (MAC) concepts to LLM agent privilege escalation \cite{ji2026seagent} and to programmable privilege control for tool-calling agents \cite{shi2025progent}.
Both represent application-layer MAC enforced by a modified agent runtime rather than by the OS kernel.
\texttt{linux-mcp} differs in that its enforcement boundary sits in the kernel itself---relying on \texttt{CAP\_NET\_ADMIN}-protected Netlink access rather than on runtime instrumentation in the agent framework---which makes the enforcement boundary independent of the daemon's correctness. The Flask security architecture \cite{spencer1999flask} and its Linux embodiment \cite{loscocco2001selinux} establish the prior art for policy/mechanism separation at the kernel level; \texttt{linux-mcp} applies the same principle to MCP-style tool admission rather than to traditional file and process access control.

The key observation across all of these is that none addresses the combination of properties the MCP control plane requires: dynamic semantic-hash-based tool identity, session-bound approval tickets with explicit lifecycle state, and durability of that state across daemon failures.
\texttt{linux-mcp} introduces a minimal kernel module for exactly these invariants.



\subsection{Threat Model}

The above limitations suggest that the correctness of MCP-style systems depends on a small set of control-plane invariants, including identity binding, approval validity, and consistency between tool semantics and execution.
We therefore ground our design in a threat model that focuses on attacks against these invariants in a local-host setting.

The local-adversary model is appropriate for MCP-style deployments for three reasons.
First, tool services in a typical deployment run on the same host as the gateway; the attack surface that most immediately needs hardening is therefore local, not remote.
Second, a remote adversary who has already achieved code execution on the host is, by definition, a local adversary: the local threat model subsumes this class without requiring a separate treatment.
Third, the two dominant remote-oriented threats in AI agent systems---prompt injection and supply-chain compromise of manifests---are explicitly handled at other layers (Section~\ref{sec:attack-scope}) and fall outside the scope of kernel-level admission control.
The threat model therefore focuses on what kernel enforcement can actually address: a co-located process that can observe, forge, or replay control-plane messages.

Concretely, we consider a local adversary operating on the same host as the MCP gateway.
The attacker interacts with the system through exposed user-space interfaces and may craft arbitrary requests, replay previously observed messages, and exploit inconsistencies in user-space state management.
We assume that the attacker has full knowledge of the system’s protocol, implementation, and manifest structure.

We do not assume elevated privileges beyond the normal user-space boundary.
Attacks that require modification of trusted configuration inputs, direct process compromise (e.g., \texttt{ptrace}), or kernel-level control are considered out of scope.
Similarly, remote adversaries and external service compromise are not part of the evaluation.

Within this scope, we focus on four classes of control-plane failures:
(1) \emph{spoofing} of identity or session context,
(2) \emph{replay} of previously valid requests or approval decisions,
(3) \emph{substitution} of tool identity or semantics, and
(4) \emph{escalation} through bypassing admission checks in the gateway.
These attack classes and their corresponding defenses are summarized in Table~\ref{tab:threat_in_scope}.

The system does not aim to provide complete host security.
Instead, it ensures that critical control-plane invariants are enforced outside the mutable user-space execution path, making them significantly harder to bypass through gateway-local bugs or weakened mediation logic.
These attack classes directly inform the design of \texttt{linux-mcp}, and are revisited in the evaluation through targeted adversarial experiments.
